\clearpage

\hypersetup{linkcolor=black}
\tableofcontents % Prints the main table of contents
\hypersetup{linkcolor=blue}
\renewcommand\arraystretch{1}

\mainmatter % Begin numeric (1,2,3...) page numbering
\pagestyle{thesis} % Return the page headers back to the "thesis" style

\chapter{Introduction}
\label{ch:introduction}

Cybersecurity is no longer a niche technical concern. It directly affects daily life through
banking, communication, healthcare, transportation, and digital public services.
Modern attacks are targeted, persistent, and increasingly systemic because software systems
are deeply connected through cloud services, CI/CD pipelines, open-source dependencies,
and distributed teams.

This report presents a broad overview of contemporary cyber threats and then narrows the focus
to software supply chain attacks. The focus is motivated by recent incidents where compromise
of a trusted upstream component propagated silently to many downstream consumers.

\chapter{Overview and Different Types of Threats}
\label{ch:threat-overview}

\section*{Declaration of Contributions}
This chapter was prepared primarily by Adesh Palkar, with technical inputs from all Group 4 members.
The threat taxonomy and mitigation mapping were consolidated jointly during group reviews.

\section{The Modern Cyber Threat Landscape}

The most common operational threat classes discussed in this report are malware/ransomware,
phishing and social engineering, denial-of-service (DDoS), insider threats, and software supply
chain attacks. A concise summary is provided in \cref{tab:threat-overview}.

\begin{table}[htbp]
\centering
\caption{High-level threat classes and representative mitigations.}
\label{tab:threat-overview}
\begin{tabular}{p{0.26\textwidth} p{0.30\textwidth} p{0.34\textwidth}}
\toprule
\textbf{Threat Class} & \textbf{Typical Attack Vector} & \textbf{Representative Mitigations} \\
\midrule
Malware and Ransomware & Malicious payload execution for encryption, theft, or disruption & Endpoint protection, offline backups, rapid patching \\
Phishing and Social Engineering & Deceptive communication that steals credentials or installs malware & MFA, secure email gateway, user awareness \\
DDoS & Flooding network/service resources to exhaust availability & Traffic filtering, rate limiting, redundancy \\
Insider Threats & Misuse of legitimate internal access & Least privilege, monitoring, segregation of duties \\
Supply Chain Attack & Compromise of trusted package, dependency, or build path & Provenance checks, signed releases, continuous vetting \\
\bottomrule
\end{tabular}
\end{table}

\section{Threat 1: Malware and Ransomware}

Malware and ransomware are among the most financially damaging threats in current practice.
Attackers use malicious payloads to encrypt files, steal credentials, or disrupt operations,
often combining endpoint compromise with data exfiltration and extortion.

\textbf{Key mitigations:} endpoint protection, rapid patching, strict execution policies,
and secure offline backups with tested recovery workflows.

\section{Threat 2: Phishing and Social Engineering}

Phishing remains the most common entry vector for multi-stage intrusions.
Adversaries exploit human trust through spoofed emails, fake portals, and impersonation to
collect credentials or trigger malware execution.

\textbf{Key mitigations:} multi-factor authentication, secure email filtering,
and continuous awareness training with simulated phishing exercises.

\section{Threat 3: Denial of Service (DDoS)}

DDoS attacks flood services with traffic and exhaust network or application resources,
causing availability failures for legitimate users.

\textbf{Key mitigations:} rate limiting, upstream filtering, traffic scrubbing,
and redundant infrastructure with failover planning.

\section{Threat 4: Insider Threats}

Insider threats originate from authorized users whose access is misused intentionally or
accidentally. These incidents can bypass perimeter controls because the actor already operates
inside trust boundaries.

\textbf{Key mitigations:} least-privilege access, monitoring of privileged activity,
and segregation of duties for sensitive operations.

\section{Threat 5: Software Supply Chain Attacks}

Software supply chain attacks differ from direct endpoint attacks because the compromised
element is often a trusted third-party package, maintainer account, build pipeline,
or release artifact. The attack propagates through normal dependency workflows.

\begin{figure}[htbp]
\centering
\fbox{\parbox{0.86\textwidth}{
\textbf{Conceptual Supply Chain Path:}
Developer publishes package $\rightarrow$ registry distributes package
$\rightarrow$ downstream projects install dependency $\rightarrow$ compromise fans out to users.
}}
\caption{Conceptual propagation path in a software supply chain attack.}
\label{fig:supply-chain-path}
\end{figure}

\section{Scope and Focus}

The broad threat overview is narrowed to supply chain attacks for three reasons:

\begin{enumerate}[label=(\roman*)]
\item the trust boundary in dependency ecosystems is often implicit,
\item one upstream compromise can impact many organizations,
\item detection and response are difficult when malicious behavior is hidden in normal update paths.
\end{enumerate}

\chapter{Case Study}
\label{ch:casestudy}

\section*{Declaration of Contributions}
This chapter was prepared primarily by Nisarg Prajapati, with supporting review by
Adesh Palkar and Utkarsh Kumar.

\section{Case Study: XZ Utils Backdoor (2024)}

XZ Utils is a widely used compression utility in Linux ecosystems and includes
the \texttt{liblzma} library used by many software stacks.
In 2024, malicious code associated with releases 5.6.0 and 5.6.1 raised serious
concerns due to potential unauthorized access paths in environments linked with
\texttt{sshd}/OpenSSH \citep{nvd2024cve3094,redhat2024xz}.

\begin{table}[htbp]
\centering
\caption{Simplified timeline and implications of the XZ Utils incident.}
\label{tab:xz-timeline}
\begin{tabular}{p{0.22\textwidth} p{0.66\textwidth}}
\toprule
\textbf{Event} & \textbf{Observation} \\
\midrule
Upstream release activity & Malicious behavior was associated with release/build artifacts.\\
Potential integration path & Systems using affected dependency paths could inherit risk transitively.\\
Ecosystem response & Advisories and coordinated mitigation were issued to reduce spread.\\
Security lesson & A trusted low-level dependency can become a high-impact attack surface.\\
\bottomrule
\end{tabular}
\end{table}

\section{Attack Path and Security Insight}

The attack path can be interpreted as follows:

\begin{enumerate}
\item compromise is introduced at an upstream package/release stage;
\item downstream systems consume the package through normal dependency workflows;
\item malicious behavior is activated in operational contexts where trust is assumed.
\end{enumerate}

This pattern demonstrates a core insight: in supply chain attacks, \emph{implicit trust}
in dependency ecosystems can bypass traditional defensive boundaries.

\section{Why This Incident Matters}

A trusted low-level dependency became a potential ecosystem-level risk amplifier.
The case highlights how delayed detection in upstream components can expose many downstream
applications simultaneously.

\chapter{Problem Statement}
\label{ch:problem}

\section*{Declaration of Contributions}
This chapter was prepared jointly by Nisarg Prajapati and Adesh Palkar,
with final wording review by all members.

\textbf{The widespread blind trust in unverified open-source dependencies exposes organizations
to devastating software supply chain attacks.}

\section{Technical Meaning}

Technically, this problem appears because:

\begin{itemize}
\item modern software depends on deep transitive dependency chains;
\item package adoption is often based on popularity or convenience rather than verified provenance;
\item release artifact integrity and maintainer-risk checks are inconsistently enforced.
\end{itemize}

\section{Why This Matters}

A single compromised package can affect many products and organizations across sectors.
Attack impact propagates from upstream maintainers to downstream applications and then to end users.
National-level advisories reinforce the need for rapid, layered defense strategies and coordinated
response \citep{cisa2024shieldsup}.

\section{What Is Needed}

\begin{itemize}
\item stronger dependency verification across all critical components;
\item robust artifact integrity checks in release and deployment workflows;
\item better maintainer-risk monitoring and trust governance;
\item continuous supply chain security monitoring across software ecosystems.
\end{itemize}

\chapter{Related Work - NDSS Research Paper}
\label{ch:related-ndss}

\section*{Declaration of Contributions}
This chapter was prepared primarily by Utkarsh Kumar,
with editorial and technical inputs from Adesh Palkar.

\section{Research Context and Motivation}

Duan \textit{et al.} presented one of the most influential ecosystem-scale studies on package
manager attacks and proposed MALOSS (MALware detection for Open Source Software), combining
metadata analysis, static analysis, dynamic analysis, and human validation \citep{duan2021ndss}.

The paper reframes supply chain defense as a scalable ecosystem problem rather than a set of
isolated package-level incidents.

\section{MALOSS Pipeline}

\subsection{Metadata Analysis}
Metadata signals such as author history, release cadence, dependency anomalies,
and suspicious package naming patterns are used to flag high-risk candidates early.

\subsection{Static Analysis}
Source code is scanned for suspicious API usage and risky data flows,
including patterns that can indicate exfiltration or remote command execution paths.

\subsection{Dynamic Analysis}
Packages are executed in sandboxed environments to observe network behavior,
file operations, and process-level effects under realistic install/import/runtime stages.

\subsection{Human-in-the-loop Validation}
Manual triage confirms high-confidence findings, improves precision,
and feeds back into detection rule refinement.

\section{Key Findings and Practical Implications}

The NDSS study reported hundreds of malicious packages across major ecosystems,
showing that typosquatting, account compromise, and delayed takedowns are recurrent.
The core implication is that continuous, automated, and multi-layered screening is
necessary for sustainable risk reduction.

\section{Mitigation Takeaways and Limits}

\begin{itemize}
\item registry-level controls: maintainer MFA, typo detection, rapid advisories;
\item developer-level controls: dependency pinning, script review, update risk gating;
\item ecosystem-level controls: faster dependent notification and coordinated response.
\end{itemize}

Even with these controls, evasion remains possible through obfuscation,
trigger-gated behavior, and delayed payload activation.

\chapter{Related Work - Dependency Confusion}
\label{ch:related-dependency-confusion}

\section*{Declaration of Contributions}
This chapter was prepared primarily by Tanay Kapadia,
with supporting technical review by Utkarsh Kumar.

\section{Dependency Manager Context}

Modern software ecosystems rely on dependency managers and registry resolution logic.
When public and private package names overlap, resolver policies can introduce a risk
that an unintended package is selected.

\section{Dependency Confusion Mechanism}

Dependency confusion exploits naming/version resolution across public and private registries.
If resolver policies prioritize an untrusted public package, an attacker can hijack a
dependency path \citep{birsan2021dependencyconfusion}.

\section{Risk Assessment by Source Tree Matching}

Ding \textit{et al.} proposed source tree matching for systematic risk assessment,
enabling comparison of intended components with actual resolved dependencies even in
obfuscated settings \citep{ding2024dependency}.

\section{Method Summary}

\begin{enumerate}
\item obtain source and component context;
\item generate abstract syntax trees;
\item optimize trees by removing low-value nodes;
\item compare AST structures for accurate component mapping.
\end{enumerate}

\section{Why Existing Trust Models Fail}

Most ecosystems optimize for package availability and developer convenience.
As a result, they can inherit the following weaknesses:

\begin{itemize}
\item weak provenance guarantees for package origin and release artifacts;
\item delayed detection of malicious behavior hidden in build logic;
\item version-selection policies that can prioritize untrusted public packages.
\end{itemize}

\chapter{Existing Solution}
\label{ch:existing-solution}

\section*{Declaration of Contributions}
This chapter was prepared primarily by Parth Baranwal,
with technical inputs from the full group.

\section{eBPF-based Runtime Security}

Current defenses against supply chain compromise can include eBPF-based runtime telemetry and
policy enforcement at kernel level, improving visibility for process, filesystem,
and network anomalies.

\section{Practical Control Families}

\begin{enumerate}
\item \textbf{Runtime visibility and enforcement:} observe suspicious behavior close to the kernel surface.
\item \textbf{Registry-side controls:} maintainer hardening, typo detection, and package vetting.
\item \textbf{Build and release integrity:} artifact signing, provenance metadata, and reproducible build checks.
\item \textbf{Automated package analysis:} pipelines such as MALOSS for ecosystem-scale screening.
\end{enumerate}

\section{Operational Constraints}

No single control family is sufficient on its own.
Effective defense requires layered deployment and continuous policy refinement.

\chapter{Proposed Solution}
\label{ch:proposed-solution}

\section*{Declaration of Contributions}
This chapter was jointly prepared by Parth Baranwal and Utkarsh Kumar,
with review support from Tanay Kapadia.

\section{Layered Mitigation Framework}

Based on the case study and related work, we propose a layered framework that combines
prevention, detection, and response.

\begin{algorithm}[htbp]
\caption{Layered Screening for Third-party Dependencies}
\label{alg:layered-screening}
\begin{algorithmic}[1]
\Require Dependency set $D$, policy rules $P$, trusted source map $T$
\Ensure Classified set of dependencies with risk levels
\ForAll{$d \in D$}
\State Validate source and maintainer metadata against $T$
\State Run static checks on install/build scripts and sensitive APIs
\State Trigger sandboxed dynamic checks for high-risk candidates
\State Assign risk score and enforce policy action (allow, quarantine, reject)
\EndFor
\State Publish alerts and update trust rules from analyst feedback
\end{algorithmic}
\end{algorithm}

\section{Static Analysis + Symbolic Execution Direction}

To reduce missed detections from input-dependent dynamic behavior,
we recommend combining static pre-screening with selective symbolic execution for
high-risk code paths.

\section{Mitigation Limits and Practical Constraints}

Main constraints include:

\begin{itemize}
\item static analysis may miss runtime-triggered behavior or heavily obfuscated logic;
\item dynamic analysis depends on execution paths and may miss dormant payloads;
\item full manual verification does not scale to very large package ecosystems;
\item symbolic comparison workflows can become computationally expensive without good heuristics.
\end{itemize}

\chapter{Conclusion}
\label{ch:conclusion}

\section*{Declaration of Contributions}
This chapter was jointly prepared by Parth Baranwal and Adesh Palkar,
with final review by all Group 4 members.

This report examined modern cyber threats and focused on software supply chain attacks
as a high-impact, trust-centric threat model. Through the XZ Utils case study and related
research synthesis, we observed that a single compromised upstream dependency can propagate
risk to many downstream systems.

The central conclusion is that dependency trust must be continuously verified, not implicitly
assumed. Existing controls improve visibility and reduce risk, but robust defense requires
layered safeguards across package publication, dependency resolution, CI/CD pipelines,
and runtime monitoring.

\section{Future Work}

The following directions are recommended for improving practical defenses:

\begin{enumerate}
\item standardized provenance and signed artifact verification across package ecosystems;
\item stronger CI/CD policy gates for dependency updates and transitive risk checks;
\item scalable hybrid analysis (static + dynamic + symbolic) with low false-positive rates;
\item faster ecosystem-wide notification and coordinated mitigation for exposed dependents;
\item broader adoption of security-by-default package publishing and maintainer hardening.
\end{enumerate}

\chapter{References}
\label{ch:references-note}

\section*{Declaration of Contributions}
This chapter was curated by the full group to ensure citation completeness and consistency.

The references used in this report are provided in the bibliography section generated at the
end of this document.
